\section{问题的提出}

巴塞尔问题是数学史上最著名的无穷级数问题之一。它要求计算
\begin{equation}
  S=\sum_{n=1}^{\infty}\frac{1}{n^2}
   =1+\frac{1}{4}+\frac{1}{9}+\frac{1}{16}+\cdots .
  \label{eq:basel}
\end{equation}
这个问题在 1644 年由意大利数学家蒙戈利（Pietro Mengoli）提出，
后来因为伯努利兄弟的研究而广为人知，因此也被称为巴塞尔问题。
欧拉在 1735 年给出了精确答案：
\begin{equation}
  \boxed{\sum_{n=1}^{\infty}\frac{1}{n^2}=\frac{\pi^2}{6}}.
  \label{eq:answer}
\end{equation}

这个结果之所以令人惊讶，是因为左边完全由整数构成，右边却出现了圆周率 $\pi$。
它也说明整数平方的倒数和与圆的几何之间存在深刻联系。
